\begin{textsample}{NEG Dim 1 – persona\_gemini – Score -1.00 – t899\_gemini.txt}  \label{ex:f1_neg_017}
Indonesia 's \$20 billion textile industry, a \textbf{supplier} to global fashion brands, has inflicted severe environmental damage.
Since the 1990s, pollution has contaminated rivers like the Citarum and made rice fields unproductive.
The economic losses from this damage amounted to IDR11.4 trillion ( US\$838 million ) between 2004 and 2015.
In May 2016, a significant court victory brought hope for cleaner rivers.
The court suspended government decrees that had legalized wastewater discharges from three companies : PT.
Kahatex, PT.
Insan Sandang Internusa, and PT.
Five Star Textile.
The ruling was based on the fact that the permits were issued without adequate studies or \textbf{monitoring}, and it invoked the precautionary principle.
Greenpeace urges global fashion brands to fulfill their Detox commitments to \textbf{safeguard} human health and protect ecosystems.

% matched lemmas: monitoring, safeguard, supplier
\end{textsample}
